\Needspace{16\baselineskip}
\SourceLeaf{104}%

\begin{center}
BEMERKUNGEN ZUR INVERSION DES ELLIPTISCHEN INTEGRALS ERSTER GATTUNG.
\end{center}

\noindent\hspace*{\parindent}Diese Notiz, deren Entstehungszeit man vermutlich nach einer Tagebuchangabe vom 6.\ Mai 1800 »Theoriam quantitatum transcendentium $\int\frac{dx}{\BareRoot{(1-a xx)(1-\beta xx)}}$ ad summam universalitatem perduximus« wird festsetzen dürfen, gehört zu den interessantesten Entwicklungen, welche \textsc{Gauss} im Gebiete der elliptischen Funktionen angestellt hat. Im Hinblick auf Problemstellung und Bezeichnungsweise schließt sich dieselbe an die Untersuchungen »Elegantiores integralis $\int\frac{dx}{\BareRoot{(1-x^4)}}$ proprietates« (ges. Werke III p.\,404 ff.) an, welche der frühesten Periode der Beschäftigung mit den lemniskatischen Funktionen angehören. Andererseits hat sich \textsc{Gauss} 1799 mit der Inversion des allgemeinen Integrals erster Gattung $\int\frac{du}{\BareRoot{(1+\mu\mu\sin u^2)}}$ beschäftigt.

Sachlich wird die hier vorliegende Entwicklung durch den oben abgedruckten Brief von \textsc{Gauss} an \textsc{Bessel} in sehr interessanter Weise beleuchtet. \textsc{Gauss} betont daselbst wiederholt die Bedeutung der ganzen transzendenten Funktionen. Hier wird in eleganter Entwicklung die Funktion $f\varphi$, welche später von \textsc{Jacobi} mit $\sin\mathrm{am}\,\varphi$ bezeichnet wurde, als Quotient zweier Funktionen $P\varphi$ und $Q\varphi$ dargestellt, welche \textsc{Gauss} unzweifelhaft in ihrer Eigenschaft als \textls{ganze} transzendente Funktionen gekannt hat.

Es sei noch die Bemerkung gestattet, dass die Funktionen $P, Q$, welche \textsc{Gauss} hier einführt, keine anderen sind als diejenigen, welche späterhin \textsc{Weierstrass} in Anknüpfung an \textsc{Abel}s Arbeiten mit der Bezeichnung $\mathrm{Al}(\varphi)_1$ und $\mathrm{Al}(\varphi)_0$ belegte.

\hfill \textsc{Fricke}.

\bigskip

